\documentclass[11pt]{article}
\usepackage[utf8]{inputenc}
\usepackage[T1]{fontenc}
\usepackage{lmodern}
\usepackage{amsmath,amssymb,amsthm}
\usepackage{enumitem}

\title{Theorems}
\author{Author Name}
\date{\today}

\newtheorem{theorem}{Theorem}[section]
\newtheorem{lemma}[theorem]{Lemma}
\newtheorem{proposition}[theorem]{Proposition}
\newtheorem*{definition*}{Definition}

\begin{document}
\maketitle

\section[1]{Preliminaries}
\setcounter{subsection}{3}
\subsection{Open Sets, Closed Sets and Borel Sets in $\mathbb{R}$}
\begin{definition*}[Sup and Inf of a Sequence of Sets]
    Let ${\{ A_n \}}_{n=1}^{\infty}$ be a countable sequence of sets that belong to a $\sigma$-algebra $\mathcal{A}$.
    Then we define the following sets that belong to $\mathcal{A}$:
    \[
        \limsup{\{ A_n \}}_{n=1}^{\infty} = \bigcap_{k=1}^\infty \bigg[\bigcup_{n=k}^\infty A_n \bigg] \; \text{and} \;
        \liminf{\{ A_n \}}_{n=1}^{\infty} = \bigcup_{k=1}^\infty \bigg[\bigcap_{n=k}^\infty A_n \bigg]
    \]
\end{definition*}

\begin{definition*}[$\sigma$-algebra]
    Given a set $X$ a collection of sets $\mathcal{A}$ is called a $\sigma$-algebra given:
    \begin{enumerate}[label=(\roman*), leftmargin=*]
        \item $\emptyset \in \mathcal{A}$
        \item $\mathcal{A}$ is closed to complements
        \item $\mathcal{A}$ is closed to a countable union of sets
    \end{enumerate}
\end{definition*}

\begin{definition*}[Borel Sigma Algebra]
    A $\sigma$-algebra will be called a Borel $\sigma$ -algebra if it is the smallest sigma algebra, containing all the of the open sets. \\
    In this case, a Borel set will refer to any set of a Borel $\sigma$-algebra.
\end{definition*}

\begin{definition*}[$G_\delta$ set]
    Any set that is a \textbf{countable union of open sets} is called a $G_\delta$ set.
\end{definition*}

\begin{definition*}[$F_\sigma$ set]
    Any set that is a \textbf{countable union of closed sets} is called a $F_\sigma$ set.
\end{definition*}

\subsection{Sequences of Real Numbers}
\setcounter{theorem}{14}
\begin{theorem}[The Monotone Convergence Theorem]
    A monotone sequence of real numbers converges if and only if it is bounded.
\end{theorem}

\begin{theorem}[The Bolzano-Weierstrass Theorem]
    Every bounded sequence of real numbers has a convergent subsequence.
\end{theorem}

\begin{theorem}[The Cauchy Convergence Criterion for Real Sequences]
\end{theorem}

\begin{definition*}[Limit-Superior and Limit-Inferior]
    Given a sequence ${(a_n)}_n$ we define:
    \[
        \limsup a_n = \lim_{n \to \infty} [ \sup {\{a_k : k \geq n\}} ]
    \]
    and:
    \[
        \liminf a_n = \lim_{n \to \infty} [ \inf {\{a_k : k \geq n\}} ]
    \]
\end{definition*}

\stepcounter{theorem}{18}
\begin{proposition}
    
\end{proposition}
\begin{enumerate}[label= (\roman*), leftmargin=*]
    \item $\limsup a_n = \ell \in \mathbb{R}$ if and only if for all $\epsilon > 0$, $a_n >  \ell - \epsilon$ for infinitely many indices and $a_n > \ell + \epsilon$ for finitely many indices.
    \item $\limsup a_n = \infty$ iff ${(a_n)}_n$ is not bounded above.
    \item \[
                \limsup{\{a_n\}} = -\liminf{\{-a_n\}}
          \]
    \item A sequence of real numbers $(a_n)$ converges to an extended real number $a$ iff
    \[
        \limsup{\{a_n\}} = \liminf{\{a_n\}} = a
    \]
    \item If $a_n \leq b_n$ for all $n$,
    \[
        \limsup {\{a_n\}} \leq \liminf {\{b_n\}}
    \]
\end{enumerate}

\end{document}

